\noindent\textbf{Elección del procedimiento.} Para $u_\theta=\Omega R^2/r$, la disipación es proporcional a $r^{-4}$. Multiplicar por el volumen anular $2\pi r\,dr$ lleva a integrar $r^{-3}$: la fracción acumulada entre $R$ y $r$ es $1-R^2/r^2$. Aislar el radio produce $r/R=(1-f)^{-1/2}$.

\begin{enumerate}

\item La fracción debe satisfacer $0<f<1$.

\[\fitmath{\mathrm{ratio}=\frac{1}{\sqrt{1-f}}}\]

\[\fitmath{\mathrm{ratio}=\frac{1}{\sqrt{1-(0{,}9)}}\approx 3{,}162278\,}\]

\item Multiplicar por el radio físico del cilindro.

\[\fitmath{r=R\cdot \mathrm{ratio}}\]

\[\fitmath{r=(0{,}015)\cdot (3{,}16227766)\approx 0{,}04743416\,\mathrm{m}}\]

\end{enumerate}
